\section{The uniform repair induction}
\label{sec:area_law_repair_induction}

We keep the recursive family $(\mathcal D_j)_{j\ge0}$ of
Definition~\ref{def:area_law_repair_stage}, with stopping scale $S_*$.
Distances are maximum-norm distances. The square
$v+[-r,r]^2$ is the closed ball of radius $r$ about $v$.

% Source: 10-geometry.tex, replacement instructions, lines 393--417;
% geometry:repair-induction, lines 471--485 and 526--541.
% Revision: adc7f1241b42e322a6451854ab7e4b4c146bf78a.
\begin{definition}[Ancestors and descendants]
    \label{def:area_law_repair_ancestors}
    \lean{TNLean.PEPS.AreaLaw.Geometry.repairMarkParent,
        TNLean.PEPS.AreaLaw.Geometry.repairMarkAncestor,
        TNLean.PEPS.AreaLaw.Geometry.repairRegionParent,
        TNLean.PEPS.AreaLaw.Geometry.repairEarlierTerminalSquares}
    \leanok
    \uses{def:area_law_repair_stage,def:area_law_repair_children,
        def:area_law_repair_successor_regions}
    A mark occurrence $q$ of $\mathcal D_{j+1}$ is a pair of a nonterminal
    parent occurrence $p$ of $\mathcal D_j$ and an offset; write
    $\pi_j(q)=p$. For $n\ge0$ the ancestor at generation $i$ of a mark of
    $\mathcal D_{i+n}$ is
    \begin{align}
        \alpha_{i,0} & =\mathrm{id},
                     & \alpha_{i,n+1} & =\alpha_{i,n}\circ\pi_{i+n}.
        \notag
    \end{align}
    A fresh identifier of $\mathcal D_{j+1}$ is born at one parent
    occurrence of $\mathcal D_j$; retained old identifiers have no parent.
    The descendants of a mark are the marks and fresh identifiers whose
    ancestor it is. Finally, let $\mathcal E_j$ be the union of the
    closed terminal squares of the generations $i<j$.
\end{definition}

% Source: 10-geometry.tex, geometry:initial-stars, lines 332--349, and
% geometry:repair-induction, lines 471--485.
\begin{definition}[Conclusion of the uniform repair induction]
    \label{def:area_law_uniform_repair_conclusion}
    \lean{TNLean.PEPS.AreaLaw.Geometry.IsUniformRepairInduction,
        TNLean.PEPS.AreaLaw.Geometry.UniformRepairInductionStatement}
    \leanok
    \uses{def:area_law_repair_ancestors,def:area_law_repair_square,
        def:area_law_isolated_star}
    Let $a_0>0$. The recursive family satisfies the uniform repair
    induction with constant $a_0$ when, for all generations $i,j$ and
    $n\ge0$:
    \begin{enumerate}
        \item every mark $p$ of $\mathcal D_j$, of location $v_p$ and scale
              $S_p$, has an isolated star of radius $a_0S_p$ in the current
              regions of $\mathcal D_j$: inside $v_p+[-a_0S_p,a_0S_p]^2$
              the frontiers are the rays from $v_p$ in the eight axis and
              diagonal directions at which the identifier of the eight
              elementary sectors changes, each sector lies in one
              identifier, and consecutive sectors have equal labels exactly
              when they have equal identifiers;
        \item distinct marks $p\ne q$ of $\mathcal D_j$ satisfy
              $\lVert v_p-v_q\rVert\ge a_0\max(S_p,S_q)$;
        \item the star square of every mark of $\mathcal D_j$ is disjoint
              from $\mathcal E_j$;
        \item distinct marks of $\mathcal D_j$ have disjoint removal
              squares;
        \item if $x$ is a fresh identifier of $\mathcal D_{j+1}$ and
              $y\ne x$ has the same label, then the closures of their
              current regions meet only at mark locations of
              $\mathcal D_{j+1}$;
        \item every descendant mark of a mark $(v,S)$ of $\mathcal D_i$
              lies within distance $2d(S)$ of $v$, and so does every birth
              region of a descendant fresh identifier;
        \item for every mark $p$ of $\mathcal D_i$ there is $n$ such that
              no mark of $\mathcal D_{i+n}$ has ancestor $p$.
    \end{enumerate}
    The statement of the uniform repair induction asserts that there are
    $a_0>0$ and $T$ such that for every $S_*\ge T$ and every width $C_0$
    the family satisfies the uniform repair induction with constant $a_0$
    whenever the lower layer index $k_0$ is sufficiently large.
    Since there are finitely many marks, the fifth item implies that every
    positive-length interface of a replacement with a distinct identifier
    joins opposite labels.
\end{definition}

% Source: 10-geometry.tex, geometry:trim-margins, lines 380--394, and
% geometry:repair-induction, lines 471--485 with lines 487--496 and
% 516--542 of its proof.
\begin{lemma}[Metric invariants of the repair]
    \label{lem:area_law_repair_metric_invariants}
    \lean{TNLean.PEPS.AreaLaw.Geometry.RepairMetricHypotheses,
        TNLean.PEPS.AreaLaw.Geometry.RepairMetricHypotheses.le_scale,
        TNLean.PEPS.AreaLaw.Geometry.RepairMetricHypotheses.separated,
        TNLean.PEPS.AreaLaw.Geometry.RepairMetricHypotheses.disjoint_repairSquare,
        TNLean.PEPS.AreaLaw.Geometry.RepairMetricHypotheses.dist_descendant_le,
        TNLean.PEPS.AreaLaw.Geometry.RepairMetricHypotheses.birth_descendant_subset,
        TNLean.PEPS.AreaLaw.Geometry.RepairMetricHypotheses.disjoint_earlierTerminalSquares,
        TNLean.PEPS.AreaLaw.Geometry.RepairMetricHypotheses.terminates}
    \leanok
    \uses{def:area_law_repair_ancestors,thm:area_law_repair_margins}
    Let $0<a_0\le1$ and $S_*>0$, suppose $d(S)\le(a_0/1000)S$ for every
    $S\ge S_*$, that every initial mark has scale at least $S_*$, and that
    distinct initial marks satisfy the separation inequality with constant
    $a_0$. Then every mark of every generation has scale at least $S_*$,
    and items two, three, four, six and seven of
    Definition~\ref{def:area_law_uniform_repair_conclusion} hold.
\end{lemma}

\begin{proof}
    \leanok
    \uses{thm:area_law_repair_margins}
    Children of a parent of scale $S$ have scale $d(S)\ge S_*$ and lie
    within distance $d(S)$ of it; distinct offsets are at distance at
    least one. Two children of one parent are therefore at least
    $d(S)\ge a_0d(S)$ apart. For children of distinct parents of scales
    $S,T$ and $M=\max(S,T)$,
    \begin{align}
        a_0M-d(S)-d(T) & \ge\Bigl(1-\frac{2}{1000}\Bigr)a_0M
        \ge a_0\max\bigl(d(S),d(T)\bigr),
        \notag
    \end{align}
    which proves separation by induction on the generation. Disjointness
    of the removal squares follows from $d(S)+d(T)<a_0M$. Since
    $2d(S)\le S$, induction along a branch gives
    $\lVert v_q-v\rVert+S_q\le2d(S)$ for every descendant $q$ of positive
    generation of $(v,S)$; a fresh birth lies in the removal square of its
    parent. Along a branch the scales at least halve, while they stay at
    least $S_*$, so every branch terminates. Finally, let $p$ be a mark of
    generation $j$ and $t$ a terminal mark of generation $i<j$. The
    ancestor $a$ of $p$ at generation $i$ is nonterminal, so $a\ne t$, and
    the star square of $p$ lies within distance $2d(S_a)$ of $v_a$. A
    common point with the terminal square of $t$ would give
    $a_0\max(S_a,S_t)\le2d(S_a)+d(S_t)\le(3a_0/1000)\max(S_a,S_t)$.
\end{proof}

% Source: 10-geometry.tex, geometry:trim-margins, lines 380--394, and
% geometry:initial-stars, lines 336--339.
\begin{lemma}[Choice of the stopping scale]
    \label{lem:area_law_repair_metric_choice}
    \lean{TNLean.PEPS.AreaLaw.Geometry.initialRepairStage_scale_separated,
        TNLean.PEPS.AreaLaw.Geometry.exists_repairMetricHypotheses}
    \leanok
    \uses{lem:area_law_repair_metric_invariants,def:area_law_repair_stage,
        thm:area_law_uniform_separated_initial_marks}
    For every $0<a_0\le1/4$ there is $T$ such that for every $S_*\ge T$
    and every width $C_0$, the hypotheses of
    Lemma~\ref{lem:area_law_repair_metric_invariants} hold whenever $k_0$
    is sufficiently large.
\end{lemma}

\begin{proof}
    \leanok
    \uses{thm:area_law_repair_margins,
        thm:area_law_uniform_separated_initial_marks}
    The margin holds for all large $S$. Initial scales are sides
    $2^{\lfloor\zeta k\rfloor}$ of cells of layers $k\ge k_0$, hence at
    least $S_*$ for $k_0$ large, and distinct initial marks are separated
    by a quarter of the larger scale.
\end{proof}

% Source: 10-geometry.tex, geometry:repair-induction, lines 471--485.
\begin{lemma}[Uniform repair induction]
    \label{lem:area_law_uniform_repair_induction}
    \uses{def:area_law_uniform_repair_conclusion,
        lem:area_law_repair_metric_invariants,
        lem:area_law_repair_metric_choice}
    For a sufficiently small fixed choice of $a_0$ and the corresponding
    threshold $S_*$, every active mark $(v,S)$ at every generation has a
    terminal-square-free isolated star $v+[-a_0S,a_0S]^2$ with the
    properties of the initial isolated-star lemma. Distinct active marks
    satisfy the same separation inequality. The simultaneous removal
    squares are disjoint. Each replacement has opposite labels across
    every positive-length interface of distinct identifiers, and its
    same-label contacts with old remnants or peer replacements occur only
    at marks of the next generation. Every descendant mark and every
    descendant birth region of a mark $(v,S)$ lies within sup distance
    $2d(S)$ of $v$. Each branch terminates after finitely many
    generations.
\end{lemma}

\begin{proof}
    \uses{lem:area_law_repair_metric_invariants,
        lem:area_law_repair_metric_choice}
    The metric assertions are
    Lemma~\ref{lem:area_law_repair_metric_invariants}. The squares of
    radius $6d(S)$ about a parent lie in its star and avoid all other
    removal squares, so each modification sees only its parent's radial
    boundaries. After normalizing the removal square to $[-1,1]^2$, the
    new boundaries lie on the eight radial rays and the four supporting
    lines of the square, and the possible new marks are the nine points of
    $\{-1,0,1\}^2$. This finite configuration gives the isolated stars of
    the children with the same constant, and the label reversal gives the
    contact assertion.
\end{proof}

% Source: 10-geometry.tex, child marks, lines 404--409 and 506--514;
% geometry:descendant-count, lines 670--681.
\begin{lemma}[A child is determined by its parent and offset]
    \label{lem:area_law_repair_parent_offset}
    \lean{TNLean.PEPS.AreaLaw.Geometry.repairMarkOffset,
        TNLean.PEPS.AreaLaw.Geometry.repairMarkParent_repairMarkOffset_injective}
    \leanok
    \uses{def:area_law_repair_ancestors,def:area_law_repair_terminal_children}
    For a child occurrence $q$ at generation $j+1$, let
    $e_j(q)\in\{0,1,2\}^2$ index its centered offset
    $e_j(q)-(1,1)$. The map
    \begin{equation}
        q\longmapsto\bigl(\pi_j(q),e_j(q)\bigr)
        \label{eq:area_law_repair_parent_offset}
    \end{equation}
    is injective.
\end{lemma}
\begin{proof}\leanok
    At generation $j+1$, a child occurrence is indexed by a nonterminal
    parent and one of the nine centered offsets. Equality of the parent
    and offset therefore gives equality of the child occurrences.
\end{proof}

% Source: 10-geometry.tex, initial marks, lines 325--363, and child marks,
% lines 404--409 and 506--514.
\begin{lemma}[Finitely many marks at each generation]
    \label{lem:area_law_repair_finite_marks}
    \lean{TNLean.PEPS.AreaLaw.Geometry.finite_repairStage_markIndex}
    \leanok
    \uses{def:area_law_repair_stage,lem:area_law_repair_parent_offset}
    For every generation $j\ge0$, the mark family $K_j$ is finite.
\end{lemma}
\begin{proof}\leanok
    \uses{lem:area_law_repair_parent_offset}
    The initial mark family is the selected finite set $M$. If $K_j$ is
    finite, the parent and offset map of
    Lemma~\ref{lem:area_law_repair_parent_offset} injects $K_{j+1}$ into
    the finite set $K_j\times\{0,1,2\}^2$. Induction proves the claim.
\end{proof}

% Source: 10-geometry.tex, geometry:descendant-count, lines 670--681;
% child marks, lines 404--409 and 506--514.
\begin{lemma}[The number of descendant marks]
    \label{lem:area_law_repair_descendant_marks}
    \lean{TNLean.PEPS.AreaLaw.Geometry.repairMarkAncestor_parent_offset_injective,
        TNLean.PEPS.AreaLaw.Geometry.card_repairMarkAncestor_fiber_succ_le,
        TNLean.PEPS.AreaLaw.Geometry.card_repairMarkAncestor_fiber_le}
    \leanok
    \uses{def:area_law_repair_ancestors,lem:area_law_repair_parent_offset,
        lem:area_law_repair_finite_marks}
    Fix a mark $p\in K_i$. For $n\ge0$, let
    \begin{equation}
        F_n(p)=\{q\in K_{i+n}:\alpha_{i,n}(q)=p\}.
        \label{eq:area_law_repair_descendant_marks}
    \end{equation}
    These sets are finite. The parent and offset map restricts to an
    injection $F_{n+1}(p)\longrightarrow F_n(p)\times\{0,1,2\}^2$.
    Consequently,
    \begin{equation}
        |F_{n+1}(p)|\le9|F_n(p)|,
        \qquad |F_n(p)|\le9^n.
        \label{eq:area_law_repair_descendant_count}
    \end{equation}
\end{lemma}
\begin{proof}\leanok
    \uses{lem:area_law_repair_parent_offset,lem:area_law_repair_finite_marks}
    Each $F_n(p)$ is a subset of the finite mark family $K_{i+n}$.
    The parent of a descendant in $F_{n+1}(p)$ belongs to $F_n(p)$ by the
    defining recursion for the ancestor map. Its parent and offset
    determine the child, which proves the stated injection and the
    first inequality. Since $F_0(p)=\{p\}$, induction gives the second.
\end{proof}
